\section{Internal theories and essential incompleteness}\label{sec:internal-theories}

\subsection{Theories and their interpretation}

Let $\mathscr R$ be a restriction category with ranges, restriction
zeros, and binary restriction coproducts. Put
$\mathscr S=\Split(\mathscr R)$ and $\mathcal D=\Tot(\mathscr S)$.
Assume that $\mathcal D$ has the categorical syntax of
Definition~\ref{def:categorical-syntax} and the regular structure
required in Appendix~\ref{app:syntax-theories}.
Let $\Phi\in\mathcal D$ be the formula object and
$j:S\hookrightarrow\Phi$ its sentence subobject. For an axiom
subobject $i_A:A\hookrightarrow S$, suppose the inference
endofunctor $H_A:\mathcal D/\Phi\to\mathcal D/\Phi$ of
Appendix~\ref{app:first-order-derivations} has an initial algebra
with conclusion morphism $c_A:P_A\to\Phi$.
The pullback in $\mathcal D$
\[
\begin{tikzcd}
 P_{A,S}\arrow[r]\arrow[d,"c_{A,S}"'] & P_A\arrow[d,"c_A"]\\
 S\arrow[r,hook,"j"'] & \Phi
\end{tikzcd}
\]
restricts conclusions to sentences. By
Theorem~\ref{thm:first-order-theory}, the regular image
$m_A:T_{A,S}\hookrightarrow S$ of $c_{A,S}$ is the inverse image along $j$ of the least deductively closed formula
subobject containing the axioms. Assume that $m_A$ is a restriction monomorphism in
$\mathscr S$: it has a partial inverse
$r_A:S\rightharpoonup T_{A,S}$ with
$r_A m_A=1_{T_{A,S}}$ and $m_A r_A=\bar r_A$.
Lemma~\ref{lem:range-regular-image} then gives the restriction
idempotent representing the sentence theory,
\[
 t_A=m_A r_A=\ran(c_{A,S}):S\rightharpoonup S.
\]
Negation preserves the empty free-variable list, so it restricts
to a total morphism $\nu:S\to S$ in $\mathcal D$.
The restriction idempotent $\nu^*t_A$ represents refutable sentences.
The sentence theory is \emph{consistent} when
$t_A\nu^*t_A=0_S$, and \emph{complete} when
$t_A\vee\nu^*t_A=1_S$.
It is \emph{essentially incomplete} when it is consistent and, for every
deductively closed sentence subobject $U\hookrightarrow S$ in the same language,
\[
 T_{A,S}\leq U,\quad U\wedge\nu^*U=\bot
 \quad\Longrightarrow\quad U\vee\nu^*U\ne\top.
\]

Let $F:\mathscr R\to\mathcal A$ be a range functor, where
$\mathcal A$ is a split restriction category. For each restriction
idempotent $e$ on an object $X$ of $\mathscr R$, choose a splitting
\[
 F(e)=m_e r_e,\qquad
 m_e:E_e\to FX,\qquad r_e:FX\rightharpoonup E_e,\qquad
 r_e m_e=1_{E_e},
\]
with $m_e$ total. Choose the identity splitting when $e=1_X$.
These choices define a restriction functor
\[
 \widehat F:\mathscr S\longrightarrow\mathcal A,\qquad
 \widehat F(X,e)=E_e,\qquad
 \widehat F(f)=r_d F(f)m_e
\]
for $f:(X,e)\rightharpoonup(Y,d)$. Its restriction to
$\mathscr R$, identified with the objects $(X,1_X)$, is $F$.
Indeed, the equations $f=dfe$ and $r_e m_e=1_{E_e}$ verify
identities and composition, while the restriction axioms give
$\overline{r_d F(f)m_e}=r_e F(\bar f)m_e$.

This extension preserves ranges. To see this, write
$a=\ran(f)\le d$, where the same restriction idempotent represents
the range in $\mathscr R$ and in $\mathscr S$.
The restriction idempotent $r_d F(a)m_d$ receives
$\widehat F(f)$. If a restriction idempotent $q$ on $E_d$
receives $\widehat F(f)$, then $m_d q r_d$ is a restriction
idempotent on $FY$ receiving $F(f)$. Consequently
$F(a)=\ran(F(f))\le m_d q r_d$, and hence
$r_d F(a)m_d\le q$. Thus
$\ran(\widehat F(f))=\widehat F(\ran(f))$.
The extension also preserves restriction zeros.

The coproduct of $(X,e)$ and $(Y,d)$ in $\mathscr S$ is
$(X+Y,e+d)$. The canonical morphism
\[
 \kappa_{e,d}=r_{e+d}\chi_{X,Y}(m_e+m_d):
 E_e+E_d\longrightarrow E_{e+d}
\]
has inverse
$(r_e+r_d)\chi_{X,Y}^{-1}m_{e+d}$. The inverse equations follow
from $m_e r_e=F(e)$, $m_d r_d=F(d)$, and
$F(e+d)=\chi_{X,Y}(F(e)+F(d))\chi_{X,Y}^{-1}$.
Thus $\widehat F$ preserves binary restriction coproducts.
Since $a\vee b=\ran[a,b]$ for restriction idempotents $a,b$ on
the same object, preservation of coproducts and ranges gives
$\widehat F(a\vee b)=\widehat F(a)\vee\widehat F(b)$.

\begin{proposition}[Interpretation of a sentence theory]\label{prop:interpreted-theory}
Let $F:\mathscr R\to\mathcal A$ be a range functor with split
target, let $\widehat F:\Split(\mathscr R)\to\mathcal A$ be the
extension by chosen splittings, and let $t_A=m_A r_A=\ran(c_{A,S})$
be a sentence theory constructed above. Then
\begin{gather*}
 \widehat F(t_A)=\ran(\widehat F(c_{A,S})),\\
 \widehat F(m_A)\widehat F(r_A)=\widehat F(t_A),\qquad
 \widehat F(r_A)\widehat F(m_A)=1_{\widehat F(T_{A,S})}.
\end{gather*}
Suppose moreover that $\mathcal A=\Par(\mathcal E,\Mono)$ for a
regular category $\mathcal E$. If the total morphism
$\widehat F(c_{A,S})$ factors in $\mathcal E$ through a
monomorphism $v:V\hookrightarrow\widehat F S$, then
$\widehat F(m_A)$ is a subobject of $v$. Put
\[
 d_V=[\widehat F S\xleftarrow{v}V\xrightarrow{v}\widehat F S]:
 \widehat F S\rightharpoonup\widehat F S.
\]
If $d_V(\widehat F\nu)^*d_V=0$, then
$\widehat F(t_A)(\widehat F\nu)^*\widehat F(t_A)=0$.
\end{proposition}
\begin{proof}
Range preservation gives the first equation, and functoriality
gives the two splitting equations. In $\Par(\mathcal E,\Mono)$,
every monomorphism of $\mathcal E$ has a partial inverse.
Lemma~\ref{lem:range-regular-image} therefore identifies the range
of $\widehat F(c_{A,S})$ with the partial identity on its regular
image in $\mathcal E$. Its splitting by $\widehat F(m_A)$ and
$\widehat F(r_A)$ identifies $\widehat F(m_A)$ with that image.
Factorization through $v$ implies
$\widehat F(t_A)\le d_V$. Monotonicity of inverse image and meet
then gives
\[
 \widehat F(t_A)(\widehat F\nu)^*\widehat F(t_A)
 \le d_V(\widehat F\nu)^*d_V=0.
\]
\end{proof}

For an interpretation with validity subobject
$v:V\hookrightarrow\widehat F S$, soundness means that there is a morphism
$k:\widehat F(P_{A,S})\to V$ in $\mathcal E$ such that
$v k=\widehat F(c_{A,S})$. If validity and its inverse image under
negation are disjoint, the proposition then gives consistency of the
interpreted sentence theory.

\subsection{Incompleteness from semantic consistency}

\begin{theorem}\label{thm:logical-comparison}
Let $\mathscr R,\mathcal A,F,U,\tau,b,a_0,a_1$ satisfy the
categorical hypotheses and definitions of
Theorem~\ref{thm:comparison}. Suppose that $\mathcal A$ is split,
that inverse images along total morphisms in $\mathcal A$ preserve binary joins
of restriction idempotents, and that $F(1_U)\ne0$.
Let $t_A=\ran(c_{A,S})$ represent a sentence theory in
$\Split(\mathscr R)$ as above, with total negation
$\nu:S\to S$, and let $\theta:U\to S$ be total in
$\Split(\mathscr R)$. Put
\[
 t=\widehat F(t_A),\qquad
 P=(\widehat F\theta)^*t,\qquad
 N=(\widehat F\theta)^*(\widehat F\nu)^*t.
\]
Assume that $t(\widehat F\nu)^*t=0$, $F(a_0)\le P$, and
$F(a_1)\le N$.
Suppose there are partial morphisms
$s:FU\rightharpoonup FU+FU$ and
$\ell:U\rightharpoonup U+U$ satisfying
\[
 s^*(1_{FU}+0_{FU})=N,\qquad
 s^*(0_{FU}+1_{FU})=P,\qquad
 F\ell=\chi_{U,U}s.
\]
Then there is a total $g:U\to U$ in $\mathscr R$ such that
\[
 (\widehat F(\theta g))^*t=0,\qquad
 (\widehat F(\theta g))^*(\widehat F\nu)^*t=0.
\]
In particular,
\[
 t\vee(\widehat F\nu)^*t\ne1_{\widehat F S},\qquad
 t_A\vee\nu^*t_A\ne1_S.
\]
\end{theorem}
\begin{proof}
Preservation of meets and zero by inverse image gives
$PN=(\widehat F\theta)^*\bigl(t(\widehat F\nu)^*t\bigr)=0$.
Theorem~\ref{thm:comparison} therefore supplies a total $g:U\to U$
with $(Fg)^*P=(Fg)^*N=0$. Since $\widehat F$ agrees with $F$ on
$\mathscr R$, composition of inverse images gives the two equalities.

If $t\vee(\widehat F\nu)^*t=1_{\widehat F S}$, inverse image along
the total morphism $\widehat F(\theta g):FU\to\widehat F S$ would
give
\[
 1_{FU}
 =(\widehat F(\theta g))^*\bigl(t\vee(\widehat F\nu)^*t\bigr)
 =0\vee0=0,
\]
contrary to $F(1_U)\ne0$.
If $t_A\vee\nu^*t_A=1_S$, applying $\widehat F$ and using
preservation of joins and inverse images would give
$t\vee(\widehat F\nu)^*t=1_{\widehat F S}$, a contradiction.
\end{proof}

\subsection{Inseparability in the source}

\begin{corollary}\label{cor:domain-incompleteness}
Let $\mathcal C$ be a restriction category with restriction zeros and
binary joins of restriction idempotents preserved by inverse image along total morphisms.
Let $a,b$ be inseparable restriction idempotents on $U$. Suppose that
$\theta:U\to S$ and $\nu:S\to S$ are total, and that a restriction
idempotent $t_0$ on $S$ satisfies $a\le\theta^*t_0$ and
$b\le\theta^*\nu^*t_0$.
For every restriction idempotent $t\ge t_0$ on $S$ with $t\,\nu^*t=0_S$,
the pair $t,\nu^*t$ is inseparable. In particular,
$t\vee\nu^*t\ne1_S$.
\end{corollary}
\begin{proof}
Suppose complementary restriction idempotents $e,e'$ on $S$ satisfy
$t\le e$ and $\nu^*t\le e'$. Put
$c=\theta^*e$ and $c'=\theta^*e'$. Inverse image preserves
meets and zero, so $cc'=\theta^*(ee')=0_U$.
Preservation of joins and totality of $\theta$ give
$c\vee c'=\theta^*(e\vee e')=\theta^*1_S=1_U$.
Monotonicity gives
$a\le\theta^*t_0\le\theta^*t\le c$ and
$b\le\theta^*\nu^*t_0\le\theta^*\nu^*t\le c'$.
Thus $c,c'$ would separate $a,b$, a contradiction.
If $t\vee\nu^*t=1_S$, the disjoint pair $t,\nu^*t$ would
itself be such a complementary separator.
\end{proof}

Visser proves essential undecidability when a recursive function maps
two recursively inseparable sets into the provable and the refutable
sentences, respectively \cite[Lemma~23]{VisserEHU}.
In synthetic computability, Kirst and Peters use a reduction of a
diagonal pair to provability and refutability to find an independent
sentence in each consistent extension with semidecidable provability
\cite[Definitions~7, 17 and Theorem~20]{KirstPeters}.
Corollary~\ref{cor:domain-incompleteness} gives inseparability of
the theorem and refutation domains. In the split case, let
$j:J\hookrightarrow S$ represent their join $t\vee\nu^*t$.
Properness means that $1_S:S\to S$ does not factor through $j$.
If the total category has a terminal object $1$, a global point
outside the join is a morphism $s:1\to S$ that does not factor
through $j$. A proper subobject can contain every global element. For
example, let a nontrivial group $G$ act on itself by left
multiplication. The inclusion $1\hookrightarrow 1+G$ of $G$-sets is
proper, but every morphism $1\to 1+G$ factors through it because the
second summand has no fixed point.

An inseparable pair on an object of $\mathscr R$ remains inseparable
in $\Split(\mathscr R)$: the fully faithful inclusion
$U\mapsto(U,1_U)$ changes neither its endomorphisms nor its
restriction idempotents. Thus
Corollary~\ref{cor:domain-incompleteness} applies to sentence
objects in $\Split(\mathscr R)$ even when the inseparable pair was
constructed in $\mathscr R$.

Section~\ref{sec:arithmetic} applies
Theorem~\ref{thm:logical-comparison} to a soundness factorization
and the lift in $S_{\PA}$ constructed in
Lemma~\ref{lem:arithmetic-proof-search}.
The numerical, Miller-sequence, and cosheaf realizations apply
Corollary~\ref{cor:domain-incompleteness}. In each case,
Appendix~\ref{app:syntax-theories} supplies the generated theorem
subobjects, and axiom inclusion gives inclusion of their ranges.

